% appendix-capture-guide.tex
% Panduan pengambilan gambar dan bukti pengujian — I2S TX Peripheral
% Tambahkan ke main.tex: \chapterappendixadd{contents/appendix/appendix-capture-guide}

\section{Panduan Pengambilan Gambar dan Bukti Pengujian}

Lampiran ini berisi panduan langkah demi langkah untuk memperoleh
gambar-gambar yang digunakan sebagai bukti pengujian dalam Bab~IV.

% ------------------------------------------------------------------
\subsection{Capture Clocking Wizard: Requested vs.\ Actual Frequency}
% ------------------------------------------------------------------

\begin{enumerate}
  \item Buka proyek Vivado dan double-click IP \textit{Clocking Wizard} pada Block Design.
  \item Masuk ke tab \textit{Output Clocks}.
  \item Pastikan kolom \textbf{Requested} dan \textbf{Actual} terlihat dengan jelas pada baris \texttt{clk\_out1} (\texttt{audio\_clk}).
  \item Ambil screenshot seluruh jendela (Windows: Win+Shift+S; Linux: \texttt{scrot}).
  \item Simpan dengan nama: \texttt{contents/chapter-4/fig-clk-wiz-actual.png}.
  \item Pastikan angka \texttt{Actual = 24.56897~MHz} terbaca jelas.
\end{enumerate}

% ------------------------------------------------------------------
\subsection{Capture Vivado Block Design}
% ------------------------------------------------------------------

\begin{enumerate}
  \item Di Vivado, buka file \texttt{*.bd} (Block Design).
  \item Atur zoom agar semua blok terlihat (MicroBlaze, AXI Interconnect, I2S TX IP, Clocking Wizard).
  \item Gunakan \textit{File} $\rightarrow$ \textit{Export} $\rightarrow$ \textit{Export Block Diagram} jika tersedia,
        atau ambil screenshot.
  \item Simpan sebagai: \texttt{contents/chapter-4/fig-vivado-bd.png}.
\end{enumerate}

% ------------------------------------------------------------------
\subsection{Capture Waveform Simulasi RTL}
% ------------------------------------------------------------------

\begin{enumerate}
  \item Jalankan simulasi testbench di Vivado Simulator (atau ModelSim).
  \item Tambahkan sinyal \texttt{i2s\_bclk}, \texttt{i2s\_ws}, \texttt{i2s\_data}, \texttt{fifo\_level}, \texttt{irq\_out}
        ke jendela waveform.
  \item Zoom pada 2--3 frame I2S lengkap (minimal 128 siklus BCLK).
  \item Verifikasi: WS berganti setiap 32 BCLK; DATA muncul dengan 1-bit delay setelah tepi WS.
  \item Ambil screenshot dan simpan sebagai: \texttt{contents/chapter-4/fig-sim-waveform.png}.
\end{enumerate}

% ------------------------------------------------------------------
\subsection{Capture ILA Waveform (Hardware-in-Loop)}
% ------------------------------------------------------------------

\begin{enumerate}
  \item Tambahkan IP \textit{Integrated Logic Analyzer} (ILA) ke Block Design.
  \item Hubungkan probe ke sinyal:
    \begin{itemize}
      \item \texttt{Probe~0}: \texttt{i2s\_bclk} (1-bit)
      \item \texttt{Probe~1}: \texttt{i2s\_ws} (1-bit)
      \item \texttt{Probe~2}: \texttt{i2s\_data} (1-bit)
      \item \texttt{Probe~3}: \texttt{fifo\_level[5:0]} (6-bit)
      \item \texttt{Probe~4}: \texttt{underrun} (1-bit)
      \item \texttt{Probe~5}: \texttt{irq\_out} (1-bit)
    \end{itemize}
  \item Generate bitstream ulang, program FPGA.
  \item Di \textit{Hardware Manager}: buka ILA core, set trigger pada tepi turun \texttt{WS}.
  \item Jalankan program Vitis, tekan \textit{Run trigger}.
  \item Setelah waveform stabil:
    \begin{itemize}
      \item Screenshot mode 48~kHz $\rightarrow$ \texttt{contents/chapter-4/fig-ila-48k.png}.
      \item Ubah \texttt{CONTROL[5]} (FS\_MODE) ke 1 melalui program Vitis, screenshot ulang
            $\rightarrow$ \texttt{contents/chapter-4/fig-ila-96k.png}.
    \end{itemize}
  \item Catat periode BCLK dan WS dari kursor ILA untuk mengisi Tabel~\ref{tab:meas} pada Bab~IV.
\end{enumerate}

% ------------------------------------------------------------------
\subsection{Capture Bukti Output Audio (PCM5102A)}
% ------------------------------------------------------------------

\begin{enumerate}
  \item Hubungkan output analog PCM5102A (\texttt{OUTL}/\texttt{OUTR}) ke osiloskop atau kartu suara PC.
  \item Jalankan program Vitis dengan frekuensi uji $\approx$1~kHz
        (\texttt{phase\_step = 11} untuk mode 48~kHz, lihat Tabel frekuensi di Bab~IV).
  \item \textbf{Opsi A — Osiloskop}:
    \begin{itemize}
      \item Atur trigger pada sinyal analog.
      \item Tampilkan 2--5 periode gelombang sinus.
      \item Ambil foto layar osiloskop atau gunakan fitur \textit{screen capture} osiloskop.
      \item Simpan sebagai: \texttt{contents/chapter-4/fig-scope-sine.png}.
    \end{itemize}
  \item \textbf{Opsi B — Audacity} (jika menggunakan line-in PC):
    \begin{itemize}
      \item Rekam beberapa detik audio masuk.
      \item Gunakan \textit{Analyze} $\rightarrow$ \textit{Plot Spectrum} untuk tampilkan FFT.
      \item Screenshot tampilan FFT (puncak harus muncul di sekitar 1~kHz).
      \item Simpan sebagai: \texttt{contents/chapter-4/fig-audacity-fft.png}.
    \end{itemize}
\end{enumerate}

% ------------------------------------------------------------------
\subsection{Catatan Format dan Kualitas Gambar}
% ------------------------------------------------------------------

\begin{itemize}
  \item Gunakan format \texttt{PNG} untuk screenshot UI, waveform, dan foto.
  \item Gunakan format \texttt{PDF} untuk diagram vektor (\textit{block diagram}, \textit{flowchart}).
  \item Resolusi minimal: lebar 1200 piksel untuk screenshot.
  \item Nama file harus konsisten dengan yang dipanggil di perintah \verb|\includegraphics{}| pada file \texttt{.tex}.
  \item Jangan ubah nama file setelah referensi sudah ada di \texttt{.tex}; jika perlu ubah, cari-ganti semua referensi terkait.
\end{itemize}
